\section{Conclusion}
% claim: C01
This controlled CASPER study found a small descriptive benefit from adding correction discrepancy to an empirical motion observer's residual: F1 increased from \RFone\ to \RCFone\ with slightly fewer normal alarms. Absolute detection remained weak. The correction variant alarmed on only \RCRecallPct\% of positive observations and produced \RCEpisodeRate\ normal episodes/hour, despite fresh alarms in all eight long scheduled intervals. Ordinary controls matched or exceeded several outcomes with different false-alarm costs.

Both primary variants retained nominal progression, so the comparison does not establish progression's isolated value. The evidence supports a narrow score-comparison finding and a caution about interpreting event hits without point coverage and normal burden. It establishes neither a novel digital twin nor general cybersecurity performance. Further design changes prompted by these exposed outcomes are exploratory; broader confirmation requires appropriately reserved independent recordings and a source-specific protocol.
